\documentclass[10pt]{article}

\usepackage[utf8]{inputenc}
\usepackage[spanish]{babel}
\usepackage{amsmath}
\usepackage[margin=2.5cm]{geometry}
\usepackage{setspace}
\usepackage{graphicx}
\usepackage{tikz}
\usetikzlibrary{circuits.logic.US}

\usepackage{booktabs}   
\usepackage{ragged2e}
\usepackage[table]{xcolor} 
\usepackage{tabularx}  
\usepackage{longtable} 
\usepackage{caption}    
\usepackage{float}
\usepackage{makecell}

\usepackage{booktabs}
\usepackage{listings}

\lstdefinestyle{rstyle}{
  language=R,
  basicstyle=\ttfamily\small,
  keywordstyle=\bfseries,
  commentstyle=\itshape,
  breaklines=true,
  frame=single,
  columns=flexible,
  keepspaces=true,
  showstringspaces=false
}

\usepackage{enumitem}
\newlist{glosario}{description}{1}
\setlist[glosario]{
    font=\normalfont\bfseries,
    leftmargin=0pt,
    labelsep=0.5em,
    itemsep=0.5em
}

\newcommand{\relleno}[1][\textbf{[COMPLETAR AQUÍ]}]{\textcolor{red}{#1}}
\usepackage{colortbl}
\usepackage[colorlinks=true,linkcolor=black,urlcolor=black,citecolor=black]{hyperref}

\definecolor{tablebackground}{rgb}{0.8, 0, 0}
\newcolumntype{x}[1]{>{\centering\arraybackslash\hspace{0pt}}p{#1}}
\newcolumntype{M}[1]{>{\centering\arraybackslash}m{#1}}
\newcolumntype{N}{@{}m{0pt}@{}}
\newcolumntype{P}[1]{>{\RaggedRight\arraybackslash}p{#1}}

\newcommand{\itemEmail}{ccapatinta@unsa.edu.pe}
\newcommand{\itemStudent}{Camilo B.A. Capatinta Almanza}
\newcommand{\itemCourse}{Lab. Estadística matemática, probabilidades y métodos empíricos}
\newcommand{\itemCourseCode}{2502223}
\newcommand{\itemSemester}{IV}
\newcommand{\itemUniversity}{Universidad Nacional de San Agustín de Arequipa}
\newcommand{\itemFaculty}{Facultad de Ingeniería de Producción y Servicios}
\newcommand{\itemDepartment}{Departamento Académico de Ingeniería de Sistemas e Informática}
\newcommand{\itemSchool}{Escuela Profesional de Ingeniería de Sistemas}
\newcommand{\itemAcademic}{2026 - B}
\newcommand{\itemInput}{06/11/2026}
\newcommand{\itemOutput}{08/11/2026}
\newcommand{\itemPracticeNumber}{03}
\newcommand{\itemTheme}{Distribuciones de probabilidad en R: hipergeométrica, binomial, Poisson y normal, con aproximaciones y gráficos Q-Q}

\usepackage{fancyhdr}
\pagestyle{fancy}
\fancyhf{}
\setlength{\headheight}{40.51407pt}
\renewcommand{\headrulewidth}{1pt}
\renewcommand{\footrulewidth}{1pt}

\fancyhead[L]{\raisebox{-0.2\height}{\includegraphics[width=3cm]{img/logo_episunsa.png}}}
\fancyhead[C]{\fontsize{7}{7}\selectfont    \itemUniversity \\ \itemFaculty \\ \itemDepartment \\ \itemSchool \\ \textbf{\itemCourse}}
\fancyhead[R]{\raisebox{-0.2\height}{\includegraphics[width=1.2cm]{img/logo_abet.png}}}

\fancyfoot[L]{Capatinta Camilo}
\fancyfoot[C]{\itemCourse}
\fancyfoot[R]{Página \thepage}

\begin{document}

	\vspace*{10px}
	
	\begin{center}
		\parbox{0.8\textwidth}{\centering\fontsize{17}{20}\selectfont\textbf{Informe de laboratorio \itemPracticeNumber}}
	\end{center}
	\begin{center}
		\parbox{0.8\textwidth}{\centering\textbf{\Large Tema: \itemTheme}}
	\end{center}

\begin{flushright}
    \begin{tabular}{|M{2.5cm}|N|}
        \hline 
        \rowcolor{tablebackground}
        \color{white} \textbf{Nota}  \\
        \hline 
             \\[30pt]
        \hline          
    \end{tabular}
\end{flushright}    

\begin{table}[h!]
    \begin{tabular}{|x{4.7cm}|x{4.8cm}|x{4.8cm}|}        
        \hline 
        \rowcolor{tablebackground}
        \color{white} \textbf{Estudiante} & \color{white}\textbf{Escuela}  & \color{white}\textbf{Asignatura}   \\
        \hline 
        {\itemStudent \par \itemEmail} & \itemSchool & {\itemCourse \par Semestre: \itemSemester \par Código: \itemCourseCode}     \\
        \hline          
    \end{tabular}
\end{table}     

\begin{table}[H]
    \begin{tabular}{|x{4.7cm}|x{4.8cm}|x{4.8cm}|}
        \hline 
        \rowcolor{tablebackground}
        \color{white}\textbf{Laboratorio} & \color{white}\textbf{Tema}  & \color{white}\textbf{Duración}   \\
        \hline 
        \itemPracticeNumber & \itemTheme & 02 horas   \\
        \hline 
    \end{tabular}
\end{table}

\begin{table}[H]
    \begin{tabular}{|x{4.7cm}|x{4.8cm}|x{4.8cm}|}
        \hline 
        \rowcolor{tablebackground}
        \color{white}\textbf{Semestre académico} & \color{white}\textbf{Fecha de inicio}  & \color{white}\textbf{Fecha de entrega}   \\
        \hline 
        \itemAcademic & \itemInput &  \itemOutput  \\
        \hline 
    \end{tabular}
\end{table}

\newpage
\tableofcontents
\thispagestyle{plain}
\clearpage

\section{Actividades pag. 146}

\subsection{Ejercicio 14. Captura y recaptura de peces: distribución hipergeométrica y su aproximación}

En un estudio sobre la población de peces de un lago, los científicos capturan peces del lago, los marcan y los liberan.

Supongamos que, durante un período de cinco días, se marcan y liberan 200 peces de cierto tipo. Como parte del mismo estudio, tres días después se capturan 20 peces de ese tipo. Sea $X$ el número de peces marcados entre los 20 capturados. Supongamos que se sabe que el lago tiene 1000 peces de este tipo en particular.

\begin{enumerate}[label=(\alph*)]
    \item La distribución de la variable aleatoria $X$ es (elija una):
    \begin{enumerate}[label=(\roman*)]
        \item binomial
        \item hipergeométrica
        \item binomial negativa
        \item Poisson
    \end{enumerate}

    \textbf{Distribución de $X$: (ii) hipergeométrica.}

        Se extraen 20 peces \emph{sin reemplazo} de una población finita de 1000, formada por dos tipos de elementos (200 marcados y 800 no marcados), y $X$ cuenta los marcados en la muestra. Como cada extracción modifica la composición de la población, los ensayos no son independientes y la probabilidad de éxito no es constante. Este es precisamente el modelo hipergeométrico:
        \[
        X \sim \text{Hipergeométrica}(M_1 = 200,\; M_2 = 800,\; n = 20).
        \]

    \item Utilice comandos de R para hallar $P(X \leq 4)$.

    Si $X \sim \text{Hipergeométrica}(M_1, M_2, n)$, la función de probabilidad de $X$ es
    \[
    P(X = x) = \frac{\dbinom{M_1}{x}\dbinom{M_2}{n-x}}{\dbinom{M_1 + M_2}{n}},
    \]
    donde $M_1$ es el número de elementos de interés (peces marcados), $M_2$ el número de elementos restantes (peces no marcados), $n$ el tamaño de la muestra y $x$ el número de elementos de interés observados en la muestra. La probabilidad acumulada se obtiene sumando los términos desde $x = 0$ hasta el valor requerido:
    \[
    P(X \leq 4) = \sum_{x=0}^{4} \frac{\dbinom{M_1}{x}\dbinom{M_2}{n-x}}{\dbinom{M_1 + M_2}{n}}.
    \]
    Con $M_1 = 200$, $M_2 = 800$ y $n = 20$:
    \[
    P(X \leq 4) = \sum_{x=0}^{4} \frac{\dbinom{200}{x}\dbinom{800}{20-x}}{\dbinom{1000}{20}} \approx 0.6301.
    \]
    Esta suma es la que calcula el comando \texttt{phyper(4, m = 200, n = 800, k = 20)} de R.

        \begin{figure}[H]
        \centering
        \includegraphics[width=0.80\textwidth]{img/1.png}
        \caption{$P(X \leq 4)$ en Rj Rditor, Jamovi}
        \end{figure}

    \item ¿Cuál de las distribuciones enumeradas en el inciso (a) puede utilizarse como aproximación a la distribución de $X$?
    
    \textbf{Distribución aproximada: (i) binomial.}

        La muestra ($n = 20$) representa solo el 2\,\% de la población ($N = 1000$). Con una fracción de muestreo tan pequeña, extraer sin reemplazo equivale prácticamente a extraer con reemplazo: la probabilidad de obtener un pez marcado cambia muy poco entre extracciones y los ensayos son casi independientes. Por lo tanto:
        \[
        X \approx \text{Binomial}\left(n = 20,\; p = \tfrac{200}{1000} = 0.2\right).
        \]

        Si $X \approx \text{Binomial}(n, p)$, entonces
        \[
        P(X = x) = \binom{n}{x}\,p^{x}(1-p)^{n-x}.
        \]
        \[
        P(X \leq 4) = \sum_{x=0}^{4} \binom{20}{x}(0.2)^{x}(0.8)^{20-x} \approx 0.6296,
        \]
        valor que calcula el comando \texttt{pbinom(4, size = 20, prob = 0.2)} de R.

    \item Utilizando la distribución aproximada, dé una aproximación de la probabilidad $P(X \leq 4)$ y compárela con la probabilidad exacta hallada en el inciso (b).

        \begin{figure}[H]
        \centering
        \includegraphics[width=0.80\textwidth]{img/2.png}
        \caption{Aproximación de la probabilidad $P(X \leq 4)$ y comparación con la hallada en el inciso (b)}
        \end{figure}

\end{enumerate}

\subsection{Ejercicio 15. Inspección de componentes electrónicos defectuosos}

En un envío de 10\,000 componentes electrónicos de cierto tipo, 300 son defectuosos. Supongamos que se seleccionan al azar 50 componentes para su inspección, y sea $X$ el número de componentes defectuosos encontrados.

\begin{enumerate}[label=(\alph*)]

    \item La distribución de la variable aleatoria $X$ es (elija una):
    \begin{enumerate}[label=(\roman*)]
        \item binomial
        \item hipergeométrica
        \item binomial negativa
        \item Poisson
    \end{enumerate}

        \textbf{Distribución de $X$: (ii) hipergeométrica.}

        Se seleccionan 50 componentes \emph{sin reemplazo} de un envío finito de 10\,000, formado por dos tipos de elementos (300 defectuosos y 9\,700 no defectuosos), y $X$ cuenta los defectuosos en la muestra. Como cada extracción modifica la composición del envío, los ensayos no son independientes y la probabilidad de obtener un componente defectuoso no es constante. Este es precisamente el modelo hipergeométrico:
        \[
        X \sim \text{Hipergeométrica}(M_1 = 300,\; M_2 = 9700,\; n = 50).
        \]

    \item Utilice comandos de R para hallar $P(X \leq 3)$.

        Si $X \sim \text{Hipergeométrica}(M_1, M_2, n)$,
        \[
        P(X = x) = \frac{\dbinom{M_1}{x}\dbinom{M_2}{n-x}}{\dbinom{M_1 + M_2}{n}}.
        \]
        Con $M_1 = 300$, $M_2 = 9700$ y $n = 50$:
        \[
        P(X \leq 3) = \sum_{x=0}^{3} \frac{\dbinom{300}{x}\dbinom{9700}{50-x}}{\dbinom{10\,000}{50}} \approx 0.9377,
        \]
        valor que calcula el comando \texttt{phyper(3, m = 300, n = 9700, k = 50)} de R.

        \begin{figure}[H]
        \centering
        \includegraphics[width=0.80\textwidth]{img/3.png}
        \caption{$P(X \leq 3)$ en Rj Rditor, Jamovi}
        \end{figure}


    \item ¿Cuál de las distribuciones enumeradas en el inciso (a) puede utilizarse como aproximación a la distribución de $X$?
        
    \textbf{Distribución aproximada: (i) binomial.}

        La muestra ($n = 50$) representa solo el 0.5\,\% del envío ($N = 10\,000$). Con una fracción de muestreo tan pequeña, extraer sin reemplazo equivale prácticamente a extraer con reemplazo: la probabilidad de obtener un componente defectuoso casi no cambia entre extracciones y los ensayos son casi independientes. Por lo tanto:
        \[
        X \approx \text{Binomial}\left(n = 50,\; p = \tfrac{300}{10\,000} = 0.03\right).
        \]

        Si $X \approx \text{Binomial}(n, p)$,
        \[
        P(X = x) = \binom{n}{x}\,p^{x}(1-p)^{n-x}.
        \]
        Con $n = 50$ y $p = 0.03$:
        \[
        P(X \leq 3) = \sum_{x=0}^{3} \binom{50}{x}(0.03)^{x}(0.97)^{50-x} \approx 0.9372,
        \]

    \item Utilizando la distribución aproximada, dé una aproximación de la probabilidad $P(X \leq 3)$ y compárela con la probabilidad exacta hallada en el inciso (b).
        
    \begin{figure}[H]
        \centering
        \includegraphics[width=0.80\textwidth]{img/4.png}
        \caption{Aproximación de la probabilidad $P(X \leq 3)$ y comparación con la hallada en el inciso (b)}
        \end{figure}

\end{enumerate}

\subsection{Ejercicio 21. Aproximaciones sucesivas de la distribución hipergeométrica}

Supongamos que se toma una muestra aleatoria simple de 200 componentes del envío de 10\,000 componentes electrónicos del Ejercicio 15, el cual contiene 300 componentes defectuosos, y sea $Y$ el número de componentes defectuosos en la muestra.

\begin{enumerate}[label=(\alph*)]
    \item La variable aleatoria $Y$ tiene una distribución hipergeométrica$(M_1, M_2, n)$, la cual puede aproximarse mediante una distribución binomial$(n, p)$, que a su vez puede aproximarse mediante una distribución Poisson$(\lambda)$. Especifique los parámetros de cada una de las distribuciones mencionadas en la última oración.

        \textbf{Parámetros de las distribuciones.}

        \begin{itemize}
            \item \textbf{Hipergeométrica:} $Y \sim \text{Hipergeométrica}(M_1 = 300,\; M_2 = 9700,\; n = 200)$, donde $M_1$ es el número de componentes defectuosos, $M_2$ el de no defectuosos y $n$ el tamaño de la muestra.

            \item \textbf{Binomial:} como la muestra es solo el 2\,\% del envío, $Y \approx \text{Binomial}(n = 200,\; p = 0.03)$, con $p = \frac{M_1}{M_1 + M_2} = \frac{300}{10\,000} = 0.03$.

            \item \textbf{Poisson:} como $n$ es grande y $p$ es pequeño, $Y \approx \text{Poisson}(\lambda = 6)$, con $\lambda = np = 200 \times 0.03 = 6$.
        \end{itemize}

    \item Utilice comandos de R para calcular la probabilidad exacta $P(Y \leq 10)$, así como las dos aproximaciones a esta probabilidad mencionadas en el inciso (a).
            
            \begin{itemize}
                \item \textbf{Hipergeométrica (exacta):}
                \[
                P(Y \leq 10) = \sum_{y=0}^{10} \frac{\dbinom{300}{y}\dbinom{9700}{200-y}}{\dbinom{10\,000}{200}} \approx 0.9615,
                \]
                calculada con \texttt{phyper(10, m = 300, n = 9700, k = 200)}.

                \item \textbf{Binomial (primera aproximación):}
                \[
                P(Y \leq 10) \approx \sum_{y=0}^{10} \binom{200}{y}(0.03)^{y}(0.97)^{200-y} \approx 0.9599,
                \]
                calculada con \texttt{pbinom(10, size = 200, prob = 0.03)}.

                \item \textbf{Poisson (segunda aproximación):}
                \[
                P(Y \leq 10) \approx \sum_{y=0}^{10} \frac{e^{-6}\,6^{y}}{y!} \approx 0.9574,
                \]
                calculada con \texttt{ppois(10, lambda = 6)}.
            \end{itemize}

               \begin{figure}[H]
                \centering
                \includegraphics[width=0.80\textwidth]{img/5.png}
                \caption{Probabilidad exacta $P(Y \leq 10)$ y aproximaciones mediante distribuciones binomial y Poisson}
                \end{figure}

\end{enumerate}

\section{Actividades pag. 157}

\subsection{Ejercicio 10. Control de calidad del espesor de neumáticos: distribución normal}

Una máquina fabrica neumáticos cuyo espesor de banda de rodamiento sigue una distribución normal con media 10 milímetros (mm) y desviación estándar 2 mm. El neumático tiene una garantía de 50\,000 millas. Para durar 50\,000 millas, el espesor de la banda de rodamiento debe ser de al menos 7.9 mm. Si se mide que el espesor de la banda de rodamiento es menor que 7.9 mm, entonces el neumático se vende como una marca alternativa con una garantía de menos de 50\,000 millas.

\begin{enumerate}[label=(\alph*)]
    \item Halle la proporción esperada de neumáticos vendidos bajo la marca alternativa.

        Si $T \sim N(\mu, \sigma^2)$, se estandariza con $Z = \dfrac{T - \mu}{\sigma} \sim N(0,1)$. Con $\mu = 10$ y $\sigma = 2$:
        \[
        P(T < 7.9) = P\left(Z < \frac{7.9 - 10}{2}\right) = P(Z < -1.05) = \Phi(-1.05) \approx 0.1469,
        \]
        calculada con \texttt{pnorm(7.9, mean = 10, sd = 2)}. Es decir, se espera que aproximadamente el 14.69\,\% de los neumáticos se venda como marca alternativa.

                \begin{figure}[H]
                \centering
                \includegraphics[width=0.80\textwidth]{img/6.png}
                \caption{Proporción esperada de neumáticos vendidos bajo la marca alternativa}
                \end{figure}

    \item La demanda de la marca alternativa de neumáticos es tal que el 30\% de la producción total debe venderse bajo el nombre de la marca alternativa. ¿A qué valor debería fijarse el espesor crítico, originalmente 7.9 mm, para satisfacer la demanda?
        
        Se busca $c$ tal que $P(T < c) = 0.30$. Sea $z_{0.30}$ el percentil 30 de la normal estándar, con $z_{0.30} \approx -0.5244$:
        \[
        c = \mu + z_{0.30}\,\sigma = 10 + (-0.5244)(2) \approx 8.9512 \text{ mm},
        \]
        calculado con \texttt{qnorm(0.30, mean = 10, sd = 2)}. El espesor crítico debe fijarse en aproximadamente 8.95 mm.

                \begin{figure}[H]
                \centering
                \includegraphics[width=0.80\textwidth]{img/7.png}
                \caption{Espesor crítico para satisfacer la demanda de la marca alternativa}
                \end{figure}

\end{enumerate}

\subsection{Ejercicio 11. Gráficos Q-Q normales con muestras simuladas}

Responda las siguientes preguntas.

\begin{enumerate}[label=(\alph*)]
    \item Utilice el comando de R \texttt{x=runif(50)} para generar una muestra simulada de tamaño 50 de la distribución uniforme$(0, 1)$ y utilice comandos como los dados en la Sección 3.5.2 para construir un gráfico Q-Q normal. ¿Podría la muestra simulada de 50 observaciones provenir de una distribución normal? Explique.

    \item Utilice el comando de R \texttt{x=rgamma(50, 1, 1)} para generar una muestra simulada de tamaño 50 de la distribución gamma$(1, 1)$ (véase el Ejercicio 13) y utilice comandos como los dados en la Sección 3.5.2 para construir un gráfico Q-Q normal. ¿Podría la muestra simulada de 50 observaciones provenir de una distribución normal? Explique.
\end{enumerate}

\end{document}
